\chapter{Судалгааны үр дүн}
% БАГШИД ҮЗҮҮЛЭХ хуулбар (2026-11-10, ажлын хувилбараас). Засварыг Chapters/Chapter4_Results.tex-д хийнэ.
\label{ch:results}

%======================================================================
% Үр дүн: 4.1 туршилт 1 (RQ1), 4.2 туршилт 2 (RQ2), 4.3 тайлбар ба өмнөх
% судалгаатай харьцуулалт (11/09: сургуулийн 5 бүлгийн бүтцээр Хэлэлцүүлгийг
% энд нэгтгэв; хязгаарлалт 5-р бүлэгт).
% Туршилт 1-ийн тоо: runtime/stt-candidate-20261030/combined-summary.json
% ба tools/bootstrap_stt_1030.py (seed 20261106, 10000 давталт).
% Туршилт 2-ын тоо: runtime/rag-w04-runs-20261104/ (papers/19-w04-locked-comparison-20261104.mn.md).
%======================================================================

\section{Туршилт 1: Дуу таних (RQ1)}
\label{sec:res-stt-all}

\subsection{Нийт үр дүн}
\label{sec:res-stt}

Хоёр систем хоёулаа хариу өгсөн 12 бичлэгийн үр дүнг \tref{tab:stt_results}-т харуулав. Эдгээрт 176 үг, 1026 тэмдэгт байна.

\begin{table}[H]
\centering
\caption{Ижил 12 бичлэгийн нийт үр дүн}
\label{tab:stt_results}
\begin{tabular}{|l|r|r|}
\hline
\textbf{Хэмжүүр} & \textbf{Whisper} & \textbf{Gemini 3.5 Transcribe} \\
\hline
Үгийн алдаа & 67 / 176 & 71 / 176 \\
WER & 38.07\% & 40.34\% \\
WER-ийн 95\%-ийн интервал & 25.0--48.7\% & 32.3--46.2\% \\
\hline
Тэмдэгтийн алдаа & 184 / 1026 & 255 / 1026 \\
CER & 17.93\% & 24.85\% \\
\hline
Нэг бичлэгийн хугацаа (медиан) & 3.98 сек & 6.96 сек \\
\hline
\end{tabular}
\end{table}

Whisper 67, Gemini 71 үг алдсан. Gemini-ийн WER 2.3 пунктээр их. Гэхдээ хоёр системийн ялгааны 95\%-ийн интервал $-11.3$-аас $+15.8$ пункт хүртэл байна. Энэ интервал тэгийг агуулж байна. Иймээс энэ 12 бичлэгээр аль систем нь үгийн алдаагаар илүү гэдгийг хэлж чадахгүй.

Тэмдэгтийн алдаа Gemini-д их гарсан: 184-ийн оронд 255 тэмдэгт. Гэхдээ энэ зөрүүний ихэнх нь он агуулсан гурван бичлэгээс гарсан. Зөв бичвэрт оныг үгээр («хоёр мянга хорин зургаан»), Gemini цифрээр бичсэн бөгөөд тоог өөрчлөхгүй дүрмээр үүнийг алдаа гэж тоолсон. Үр дүнг харсны дараа он агуулаагүй 9 бичлэгээр тусад нь тооцоход тэмдэгтийн алдаа Whisper-т 19.33\%, Gemini-д 18.33\% байв. Энэ тооцоо урьдчилж төлөвлөөгүй тул дүгнэлт биш, зөвхөн тайлбар.

\subsection{Бичлэг бүрийн үр дүн}
\label{sec:res-per-utt}

Бичлэг бүрийн үгийн алдааг \tref{tab:stt_per_utt}-т харуулав. Gemini 7 бичлэгт Whisper-ээс бага, 5 бичлэгт их алдсан. Ижил алдаатай бичлэг байхгүй.

\begin{table}[H]
\centering
\caption{Бичлэг бүрийн үгийн алдаа}
\label{tab:stt_per_utt}
\begin{tabular}{|l|r|r|r|r|r|}
\hline
\multirow{2}{*}{\textbf{Бичлэг}} & \multirow{2}{*}{\textbf{Үг}} & \multicolumn{2}{c|}{\textbf{Whisper}} & \multicolumn{2}{c|}{\textbf{Gemini}} \\
\cline{3-6}
 & & \textbf{Алдаа} & \textbf{WER} & \textbf{Алдаа} & \textbf{WER} \\
\hline
q01 & 11 & 5 & 45.5\% & 4 & 36.4\% \\
q02 & 8 & 2 & 25.0\% & 0 & 0.0\% \\
q03 & 10 & 1 & 10.0\% & 5 & 50.0\% \\
q04 & 10 & 8 & 80.0\% & 2 & 20.0\% \\
q07 & 13 & 2 & 15.4\% & 6 & 46.2\% \\
q08 & 14 & 0 & 0.0\% & 5 & 35.7\% \\
q09 & 12 & 6 & 50.0\% & 4 & 33.3\% \\
q11 & 15 & 6 & 40.0\% & 9 & 60.0\% \\
q12 & 19 & 4 & 21.1\% & 8 & 42.1\% \\
q14 & 13 & 6 & 46.2\% & 5 & 38.5\% \\
q15 & 11 & 6 & 54.5\% & 4 & 36.4\% \\
q18 & 40 & 21 & 52.5\% & 19 & 47.5\% \\
\hline
\textbf{Нийт} & \textbf{176} & \textbf{67} & \textbf{38.1\%} & \textbf{71} & \textbf{40.3\%} \\
\hline
\end{tabular}
\end{table}

Хоёр систем өөр өөр бичлэг дээр алддаг. Жишээ нь q04-т Whisper 10 үгийн 8-ыг алдсан бол Gemini 2-ыг алдсан. Харин q08-д Whisper бүгдийг зөв таньсан бол Gemini 5 үг алдсан. Хамгийн урт q18 бичлэг (40 үг) хоёр системийн нийт алдааны бараг гуравны нэгийг бүрдүүлж байна.

\subsection{Алдааны шинжилгээ}
\label{sec:res-errors}

\textbf{Он.} q03 бичлэгт хэлсэн «2026» оныг Gemini «2020» гэж таньсан. Whisper оныг зөв таньсан. Сургуулийн мэдээлэлд он, огноо чухал. Буруу он буруу хариулт руу хөтөлнө. Иймээс энэ алдаа нь нэг үгийн алдаанаас илүү ноцтой.

\textbf{Тоо.} Gemini тоог цифрээр бичих нь элбэг. Бидний дүрмээр «2026» ба үгээр бичсэн он ялгаатай гэж тооцогдоно (\ref{sec:stt-metrics}). Энэ нь Gemini-ийн WER-ийг зарим талаар өсгөсөн байж болно. Харин q03-ийн алдаа нь бичих хэлбэрийн асуудал биш. Утга нь өөр он болсон.
% ДАРАА: тоог ижилсүүлсэн дүрмээр дахин тооцож, хоёр хувилбарыг зэрэг харуулах.

\textbf{Хугацаа.} Whisper нэг бичлэгийг медианаар 3.98 секундэд, Gemini 6.96 секундэд текст болгосон. Gemini-ийн хугацаанд интернэтээр файл илгээх хугацаа орсон (\ref{sec:stt-procedure}). Иймээс энэ нь загварын хурдыг шууд харьцуулсан тоо биш, хэрэглэгчийн хүлээх хугацаа юм.

\subsection{Урьдчилсан шалгуурын дүн}
\label{sec:res-criteria}

\ref{sec:stt-criteria}-т тогтоосон гурван нөхцөлийн аль нь ч хангагдаагүй. Gemini-ийн үгийн алдаа Whisper-ээс бага биш, их гарсан. Тэмдэгтийн алдаа муудсан. Whisper зөв таньсан оныг Gemini алдсан. Иймээс энэ хэмжилтээр Gemini-г Whisper-ээс илүү гэж үзэхгүй.

%======================================================================
\section{Туршилт 2: Хайлт ба эхтэй хариулт (RQ2)}
\label{sec:res-rag}

\subsection{Хайлт}
\label{sec:res-retrieval}

Эхтэй 20 асуултын зөв хэсэг хайлтын эхний 1, 3, 5 үр дүнд хэдэн удаа орсныг \tref{tab:rag_retrieval}-т харуулав.

\begin{table}[H]
\centering
\caption{Хайлтын үр дүн, эхтэй 20 асуулт (одоогийн тохиргоо, $k=30$)}
\label{tab:rag_retrieval}
\begin{tabular}{|l|r|r|r|}
\hline
\textbf{Хайлтын арга} & \textbf{Эхний 1-д} & \textbf{Эхний 3-т} & \textbf{Эхний 5-д} \\
\hline
BM25 & 13 & 18 & 19 \\
E5 & 14 & 19 & 19 \\
RRF (BM25 + E5) & 14 & 20 & 20 \\
Үндсэн хайлт & 15 & 20 & 20 \\
\hline
\end{tabular}
\end{table}

Үндсэн хайлтад 20 асуултын бүгдийн зөв хэсэг эхний 3 үр дүнд орсон. 15 асуултад зөв хэсэг хамгийн дээр гарсан. BM25 ба E5 дангаараа үүнээс бага байсан. Хоёуланг нь нэгтгэхэд эхний 3-т орох тоо 20 болсон.

\subsection{Хариулт}
\label{sec:res-answers}

Одоогийн тохиргооны (A) хариултын үр дүнг \tref{tab:rag_answers}-т харуулав.

\begin{table}[H]
\centering
\caption{Хариултын үр дүн, одоогийн тохиргоо (A)}
\label{tab:rag_answers}
\begin{tabular}{|p{8.0cm}|r|}
\hline
\textbf{Хэмжүүр} & \textbf{Тоо} \\
\hline
Эхтэй 20: шаардсан утга бүгд гарсан & 13 \\
Эхтэй 20: утга хэсэгчлэн гарсан & 2 \\
Эхтэй 20: татгалзсан & 5 \\
\hline
Эхгүй 10: эхгүйгээр хариулсан & 4 \\
Ресепшний эхгүй 9: эхгүйгээр хариулсан & 3 \\
\hline
\end{tabular}
\end{table}

Хайлт зөв хэсгийг олсон ч хариулт 20-оос 13-т л бүрэн гарсан. Иймээс алдааны ихэнх нь хайлт дээр биш, хариулт сонгох алхам дээр гарсан. Эх нь байхгүй 19 асуултын 7-д систем татгалзах ёстой байсан ч ойролцоо хэсгээр хариулсан. Ресепшний эхтэй ганц асуултад систем татгалзсан.

\subsection{Алдааны жишээ}
\label{sec:res-rag-errors}

Хариулт алдсан шалтгааныг асуулт бүрээр уншиж ангилсан:
\begin{itemize}
    \item \textbf{Асуух үгийг салгаж чадаагүй.} «Хэддүгээр» гэдэг асуух үг «хэддүг» болж хувирсны дараа хасагдаагүй. Иймээс эхэд байхгүй үг гэж үзээд татгалзсан.
    \item \textbf{Хэллэг эхтэй таараагүй.} Асуултад «индекс код» гэсэн байхад эхэд зөвхөн «индекс» гэж бичсэн тул татгалзсан.
    \item \textbf{Хүснэгтийн мөр.} Хариулт хүснэгтийн нэг мөрөнд байсан бөгөөд мөрөнд асуултын бусад үг байгаагүй тул татгалзсан.
    \item \textbf{Утгыг андуурсан.} «Хэдэн багш оролцсон» гэсэн асуултыг тэнхимийн багшийн тоо гэж ойлгоод тодруулга асуусан.
    \item \textbf{Эхгүй асуултад хариулсан.} Эх нь байхгүй асуултад гол үг нь давхцсан урт хэсгээр хариулсан.
\end{itemize}

\subsection{Өөрчлөлтүүдийн дүн}
\label{sec:res-rag-variants}

\tref{tab:rag_variants}-т урьдчилсан дүрмээр шалгасан B, C өөрчлөлтийг одоогийн тохиргоотой харьцуулав.

\begin{table}[H]
\centering
\caption{Тохиргооны харьцуулалт}
\label{tab:rag_variants}
\begin{tabular}{|l|c|c|c|c|}
\hline
\textbf{Тохиргоо} & \textbf{Эхтэй 20:} & \textbf{Эхгүй 10:} & \textbf{Ресепшн 9:} & \textbf{Хайлт:} \\
 & \textbf{бүрэн / дутуу / татгалзсан} & \textbf{эхгүй хариулт} & \textbf{эхгүй хариулт} & \textbf{эхний 1-д} \\
\hline
A (одоогийн) & 13 / 2 / 5 & 4 & 3 & 15 \\
B (нэг хэсэг) & 8 / 1 / 11 & 2 & 3 & 15 \\
C ($k=60$) & 12 / 2 / 6 & 4 & 3 & 14 \\
\hline
\end{tabular}
\end{table}

B эхгүй асуултын хоёр хариултыг зогсоосон. Гэхдээ эхтэй асуултын зөв хариултыг 13-аас 8 болгож бууруулсан. C хайлтын эхний байрыг 15-аас 14, зөв хариултыг 13-аас 12 болгосон. Урьдчилсан дүрмээр хоёр өөрчлөлт хоёулаа тэнцээгүй.

F1--F3 засварыг хоёр дахь түгжсэн багцаар шалгасан дүнг \tref{tab:rag_f1f3}-т харуулав.

\begin{table}[H]
\centering
\caption{F1--F3 засварын дүн}
\label{tab:rag_f1f3}
\begin{tabular}{|p{5.4cm}|c|c|}
\hline
\textbf{Багц ба тохиргоо} & \textbf{Эхтэй: бүрэн / дутуу / татгалзсан} & \textbf{Эхгүй: эхгүй хариулт} \\
\hline
Хоёр дахь багц, засваргүй & 12 / 2 / 2 (16-аас) & 1 (8-аас) \\
Хоёр дахь багц, F1--F3-тай & 12 / 2 / 2 (16-аас) & 1 (8-аас) \\
\hline
40 асуулттай багц, F1--F3-тай & 14 / 4 / 2 (20-оос) & 4 (10-аас) \\
\hline
\end{tabular}
\end{table}

Хоёр дахь багц дээр засвартай ба засваргүй хувилбар асуулт бүрээр ижил хариулсан. «Дор хаяж нэг үзүүлэлт сайжрах» нөхцөл хангагдаагүй тул гурван засварыг бүгдийг буцаасан. 40 асуулттай багц дээр засвар хариулсан асуултын тоог нэмсэн ч хоёр асуултад утга дутуу хариулт гарсан. Энэ багцын алдааг харж засвар бичсэн тул тэнд гарсан ахицыг нотолгоо гэж үзэхгүй. Иймээс гурван өөрчлөлт бүгд тэнцээгүй бөгөөд одоогийн тохиргоо (A) хэвээр үлдсэн.

%======================================================================
\section{Туршилт 3: Ресепшний асуулт}
\label{sec:res-reception}

Түгжсэн 30 асуултын дүнг гурван хэмжилтээр \tref{tab:reception_runs}-т харуулав. Хэмжилт бүрийн хооронд хийсэн засварыг доор бичив.

\begin{table}[H]
\centering
\caption{Ресепшний 30 асуултын дүн (зөв / нийт)}
\label{tab:reception_runs}
\begin{tabular}{|p{5.2cm}|c|c|c|}
\hline
\textbf{Хэмжилт} & \textbf{Хүснэгт} & \textbf{Сайт} & \textbf{Хариулах ёсгүй} \\
\hline
1. Анхны (хуваарь орсны дараа) & 8/10 & 7/10 & 9/10 \\
2. Хүснэгтийн засварын дараа & 10/10 & 7/10 & 10/10 \\
3. Хэллэгийн засвар, сайтын шинэ хуулбар & 10/10 & 10/10 & 10/10 \\
\hline
\end{tabular}
\end{table}

Анхны хэмжилтэд хүснэгтийн хоёр асуулт алдсан. «Карьер хөгжлийн төв» гэсэн товч нэрийг алба гэж таниагүй, «405 өрөө хэддүгээр давхарт вэ» гэсэн асуултыг албаны нэр гэж андуурсан. Хариулах ёсгүй асуултаас «Захирлын гар утасны дугаар хэд вэ» гэсэн асуултад систем захирлын мэндчилгээний хэсгээр хариулсан. Эдгээрийг засаж, товч нэрийг таних, бүртгэлгүй нэрийг үндсэн хайлт руу дамжуулах, хүний гар утасны дугаарыг хувийн мэдээлэл гэж татгалзах дүрэм нэмсэн.

Сайтын гурван асуултад систем алдсан. «1-р курсийн сургалтын төлбөр», «1-р курсээс элсвэл хэдэн жил суралцах» гэсэн асуултыг асуултад байгаа «1» тооноос болж бэлэн хариулттай холбож чадаагүй. Харин «Сургууль хэдэн цагт ажилладаг вэ» гэсэн асуултад олимпиадын цагийг хариулсан. Энэ бол татгалзахаас илүү аюултай буруу хариулт. Ерөнхий хэллэгийг бэлэн хариулттай холбох дүрэм нэмсний дараа гурвуулаа зөв боллоо.

Гуравдугаар хэмжилтийн өмнө сайтын хуулбарыг шинэчилсэн. Хуучин хуулбар нь хоёр тэнхимийн хуучин эрхлэгчийг харуулж байсан бол шинэ хуулбар нь сургуулийн орон тооны жагсаалттай таарсан. Энэ хэмжилтээр 30 асуулт бүгд зөв гарсан. Гэхдээ засварыг энэ багцын алдааг харж хийсэн тул энэ дүн хөгжүүлэлтийн дүн юм. Ерөнхий нарийвчлалыг шинэ, хүн бичсэн асуултаар хэмжих шаардлагатай.

\textbf{Facebook чатбот.} Сургуулийн Facebook хуудасны чатбот хариултаа Meta-ийн хиймэл оюун ухаанаар үүсгэдэг байв. Номын сангийн давхрыг асуухад «батлагдсан эх сурвалжид алга» гээд асуултыг ажилтанд шилжүүлсэн бөгөөд үүний дараа хиймэл оюун ухаан хариулахаа больсон. Иймээс зөвхөн нэг асуултаар тулгах боломжтой болсон. Мөн өмнөх ярианд чатбот олимпиадын шагналтанд 30\% хүртэл хөнгөлөлттэй гэж хариулсан. Энэ мэдээлэл сургуулийн сайтын 2026.11.09-ний хуулбарт олдоогүй. Иймээс сувгууд хооронд мэдээлэл зөрж болохыг харуулж байна.

%======================================================================
\section{Үр дүнгийн тайлбар}
\label{sec:res-interpretation}
% ӨӨРИЙН ҮГЭЭР ЗАСАХ: энэ хэсэг таны өөрийн тайлбар байх ёстой.

\subsection{Дуу таних}

Whisper-ийн анхны өгүүлэлд олон хэлний Whisper large-v2 загвар Монгол хэлэнд 100\%-иас их WER авсан \cite{radford2023whisper}. Манай Монгол хэлээр нэмэлт сургасан Whisper small загвар 38\% WER авсан. Mengke нар ч Whisper-ийг Монгол өгөгдлөөр сургахад алдаа их буурсныг харуулсан \cite{mengke2025khalkha}. Энэ нь Монгол хэлэнд нэмэлт сургалт чухал гэдгийг дэмжиж байна. Гэхдээ эдгээр өгүүллийн өгөгдөл манайхаас өөр тул тоог шууд харьцуулахгүй.

Gemini 3.5 Transcribe нэмэлт сургалтгүйгээр Монгол хэлээр сургасан Whisper-тэй ойролцоо үгийн алдаатай гарсан. Тэмдэгтийн алдааны зөрүү гол төлөв оныг цифрээр бичсэнээс гарсан. Гэхдээ нэг бичлэгт оныг буруу (2026-г 2020 гэж) таньсан. Sidiropoulos нарын туршилтаар яриа танилтын алдаа хайлтын үзүүлэлтийг хагасаас илүү бууруулсан. Үгийг буруу хуваах, нэрийг буруу бичих нь хамгийн элбэг алдаа байсан \cite{sidiropoulos2022asr}. Иймээс ресепшнд он, тоо, нэрийн алдаа буруу хариулт руу хөтөлж болно. Манай систем ийм үгийг хэрэглэгчид «нягтлаарай» гэж харуулдаг нь үүнтэй холбоотой.

2026 оны 11-р сард системд Gemini 3.5 Transcribe-ийг ашиглахаар шийдсэн. Энэ нь хэмжилтийн дүнгээс гарсан шийдвэр биш гэдгийг шууд хэлэх хэрэгтэй. Урьдчилсан гурван шалгуурын аль нь ч Gemini-д хангагдаагүй, Whisper ч энгийн компьютерын CPU дээр хурдан ажилласан (\tref{tab:stt_results}). Gemini-г дараах хэрэглээний шалтгаанаар сонгосон: загварыг дахин сургахгүйгээр сургуулийн нэрсийг жагсаалтаар зааж өгөх боломжтой, ресепшний компьютерт загвар суулгаж шинэчлэх шаардлагагүй. Гэхдээ нэрийн жагсаалт нэг хүний 8 бичлэг дээр алдааг бууруулаагүй. Үнэгүй түвшинд туршилтын үеэр хязгаарын алдаа гарсан тул бодит ашиглалтад төлбөртэй түвшин хэрэгтэй. Мөн интернэт, API түлхүүр хэрэгтэй, аудио Google руу очдог. Whisper-ийн дүн харьцуулалтын суурь болж үлдэнэ. Шинэ хэлэгчийн бичлэг ирвэл энэ шийдвэрийг дахин шалгана.
% ӨӨРИЙН ҮГЭЭР ЗАСАХ: шийдвэрийн шалтгааныг багштай ярилцсан үгээрээ бич
% (багшийн үгийг зохиохгүй).

\subsection{Хайлт ба эхтэй хариулт}

Хайлт сайн, хариулт сул гэсэн дүр зураг өмнөх судалгаанд ч гарсан. Burch их сургуулийн туслах бүх асуултын зөв хуудсыг эхний 5-д олсон ч асуултын 25\%-д хариулт агуулаагүй өөр хэсгийг авсан \cite{matic2026burch}. Манай систем ч бүх эхтэй асуултын зөв хэсгийг эхний 3-т олсон ч 20-оос 13-т л бүрэн хариулсан. Манай алдааны шалтгаан өөр: хариулт сонгох дүрэм асуултын үгийг эхийн үгтэй шууд тулгадаг. Монгол хэлэнд үгийн төгсгөл олон янзаар хувирдаг тул «хэддүгээр», «индекс код» зэрэг хэллэг эхтэй таарахгүй байна. Хариултыг шинээр бичихгүй байх нь хариулт бүрийг эхтэй нь шууд холбодог. Сул тал нь хэллэг бага зэрэг өөр байхад хариулт олдохгүй байх.

E5 нь BM25-аас бага зэрэг илүү, хоёуланг нь нэгтгэхэд бүх зөв хэсэг эхний 3-т орсон. Энэ нь Cormack нарын RRF дан системээс илүү гарсан үр дүнтэй нийцэж байна \cite{cormack2009rrf}. Гэхдээ манай 20 асуулт цөөн тул аль арга илүү гэж дүгнэхгүй.

Эх нь байхгүй 19 асуултын 7-д систем ойролцоо хэсгээр хариулсан. Ресепшнд буруу хариулт «мэдэхгүй» гэхээс илүү хор хөнөөлтэй: хүн буруу өрөө рүү явж болно. SQuAD 2.0-ийн туршилтаар ч хариулт байхгүйг таних нь хариулт олохоос хэцүү байсан \cite{rajpurkar2018squad2}. B өөрчлөлт эхгүй хариултыг цөөлсөн ч зөв хариултыг ч хамт цөөлсөн. Нэг дүрмийг чангатгах нь хоёр талдаа нөлөөлж байна. Marcel-д хайлтыг төгс (зөв хариутай) болгоход эхгүй асуултад зөв татгалзах нь 0.68-аас 0.93 болсон. Зохиогчид илүү болгоомжтой хайлт татгалзалтыг сайжруулна гэж үзсэн \cite{trienes2025marcel}. Иймээс эхгүй асуултыг тусад нь танидаг аргыг судлах хэрэгтэй.

\subsection{Ресепшний асуулт}

Ресепшний 10 асуултын 9-д сайтад эх байгаагүй. Систем тэдгээрийн 3-т эхгүйгээр хариулсан. Эх нь байсан ганц асуултад ч татгалзсан. Өөрөөр хэлбэл одоогийн сантай систем ресепшний гол асуултад хариулж чадахгүй. Энэ нь хайлтын арга биш, өгөгдлийн дутагдал юм. Marcel-ийн зохиогчид ч хэрэглэгчийн асуултаас сангийн дутуу хэсгийг илрүүлсэн \cite{trienes2025marcel}. ALOHA систем хуваарь, нээлттэй цаг зэрэг хүснэгтэн мэдээллийг тусад нь хадгалж хайдаг \cite{tao2025aloha}. Манай систем ч өрөө, багш, хуваарийг тусдаа хүснэгтээр хайх хэсэгтэй (\ref{sec:system}). Сургалтын албаны мэдээлэл орж ирсний дараа ресепшний асуултыг шинэ багцаар хэмжсэн (
ef{sec:res-reception}). Хүснэгт орсны дараа өрөө, алба, ангийн асуултад систем хариулдаг болсон.
